\begin{remark}[Optimality and a convenient weaker form]
Let $d\geq1$, let $\varepsilon\in\{-1,0,1\}$, and prescribe total mass
$Dd$, first moment zero and second moment $D\varepsilon$ for a positive
measure on the unit circle. The bound in
\cite[Proposition~\ref{B-prop:optimized-moment}]{Finite} is optimal for
these three moments. Indeed,
put $T=Dd$ and choose $\theta$ with
$\cos\theta=-(d+\varepsilon)/(2d)$. The measure with mass
\[
 \frac{T(d+\varepsilon)}{3d+\varepsilon}\text{ at }1,
 \qquad
 \frac{Td}{3d+\varepsilon}\text{ at each of }e^{i\theta},e^{-i\theta}
\]
has total mass $T$, first moment zero, and second moment $D\varepsilon$,
and it attains equality. If the two latter points coincide, their masses
are added. This optimality statement does not assert that these points are
roots of unity or that the measure is categorically realizable.

For calculations, the simpler Fej\'er polynomial $1+z+z^2$ gives
\[
 9L_X\leq D(3d+2\varepsilon)\leq D(3d+2).
\]
The optimized upper bound is smaller than this last upper bound by exactly
$2D/[9(3d+1)]$. Both estimates apply without a unitary structure;
pseudo-unitarity supplies the positive dimensions used in their proofs.
\end{remark}
